\chapter{扩展讨论与跨场景对比}

\section{跨数据集迁移}
当模型迁移到不同采样频率或不同噪声结构的数据集时，原有参数先验可能失效。建议采用“先验缩放 + 小步长预热 + 分层放开参数”的迁移流程。

\section{目标函数重加权}
对于异方差观测，可采用加权最小二乘：
\[
\Phi_w(\theta)=\frac{1}{2}\sum_{i=1}^m w_i\|r_i(\theta)\|_2^2.
\]
权重可依据观测可信度或局部方差估计动态更新。

\section{计算与精度折中}
高精度参数恢复通常伴随更高计算成本。工程中可基于业务阈值设计分级策略：常规模式优先速度，校准模式优先精度。

\section{可解释性补充}
参数解释应结合物理意义与灵敏度排序。对低敏感参数，不宜过度解释其细微波动，应通过区间形式报告不确定性。

\section{本章结语}
扩展讨论表明，参数确定问题不仅是数值优化问题，也是数据质量、先验信息和工程约束的综合平衡问题。
